% !TeX root = ../main.tex
\chapter{国内外研究现状}
\label{chap:literature}

\section{相关概念与研究分类}
先界定关键概念，再按照研究问题、方法或技术路线组织文献。
下面的两条参考文献仅用于演示引用排版，与实际课题无关，撰写时应全部替换。
单篇引用示例：\LaTeX{} 文档系统的相关说明见文献\cite{lamport1994}。
多篇引用示例：文献资料可以通过统一的文献库管理\cite{lamport1994,knuth1984}。

\section{代表性工作比较}
围绕相同维度比较各类方法，说明它们解决了什么问题、采用什么假设，以及在哪些条件下仍存在局限。
不要仅按作者和年份逐篇罗列。表\ref{tab:literature}展示了比较表的写法，表中内容均为待替换提示。

\begin{table}[htbp]
  \centering
  \caption{代表性研究的比较维度示例}
  \label{tab:literature}
  \begin{tabularx}{\linewidth}{@{}p{22mm}XXX@{}}
    \toprule
    研究类别 & 核心思路 & 适用条件 & 主要局限 \\
    \midrule
    方法一 & 填写关键方法 & 填写必要假设 & 填写尚未解决的问题 \\
    方法二 & 填写不同思路 & 填写资源或数据条件 & 填写可验证的不足 \\
    \bottomrule
  \end{tabularx}
\end{table}

\section{现有工作的不足与本课题切入点}
从上述比较中归纳研究空缺，区分文献已明确指出的不足与本课题提出的判断。
说明拟开展的工作与代表性方法的关系，并自然引出下一章的研究目标。
